\documentclass[11pt,a4paper]{amsart}

\usepackage{amsmath,amssymb,amsthm}
\usepackage[colorlinks=true,linkcolor=black,citecolor=black,urlcolor=black]{hyperref}
\usepackage[margin=1in]{geometry}

\theoremstyle{plain}
\newtheorem{theorem}{Theorem}[section]
\newtheorem{proposition}[theorem]{Proposition}
\newtheorem{lemma}[theorem]{Lemma}
\newtheorem{corollary}[theorem]{Corollary}

\theoremstyle{definition}
\newtheorem{definition}[theorem]{Definition}
\newtheorem{hypothesis}[theorem]{Hypothesis}
\newtheorem{gap}[theorem]{GAP}
\newtheorem{remark}[theorem]{Remark}

\newcommand{\RH}{\mathrm{RH}}
\newcommand{\GRH}{\mathrm{GRH}}
\newcommand{\ERH}{\mathrm{ERH}}
\newcommand{\Real}{\operatorname{Re}}
\newcommand{\Imag}{\operatorname{Im}}
\newcommand{\ord}{\operatorname{ord}}
\newcommand{\dist}{\operatorname{dist}}
\newcommand{\pair}{\mathrm{pair}}
\newcommand{\comp}{\mathrm{comp}}
\newcommand{\raw}{\mathrm{raw}}
\newcommand{\val}{\mathrm{val}}
\newcommand{\corr}{\mathrm{corr}}
\newcommand{\rem}{\mathrm{rem}}
\newcommand{\bg}{\mathrm{bg}}
\newcommand{\loc}{\mathrm{loc}}
\newcommand{\calF}{\mathcal F}
\newcommand{\calC}{\mathcal C}
\newcommand{\C}{\mathbb C}
\newcommand{\R}{\mathbb R}

\title{A Paired Dirichlet Source Package and the GRH Transfer Skeleton}
\author{Working draft}
\date{April 2026}

\subjclass[2020]{11M06, 11M26, 11M35}
\keywords{Dirichlet \(L\)-functions, generalized Riemann hypothesis,
completed \(L\)-functions, source package, local transfer}

\begin{document}

\begin{abstract}
This is a proof skeleton for the GRH transfer program suggested by the current
RH manuscript and the GRH Dirichlet-tau research cycle.  The paper separates
three layers.  First, it proves a completed-Hadamard source package for
primitive nonprincipal Dirichlet pairs, in the normalization currently cleared
by the verifier lane.  Second, it records the local source-to-slot mechanism
as a conditional theorem whose missing input is the paired unit-coordinate
local normal form.  Third, it states the formal GRH-family transfer theorem
with every currently missing analytic or local step marked as a GAP.  The
zeta-specific source theorem is treated as an input from the RH manuscript,
not reproved here.
\end{abstract}

\maketitle
\tableofcontents

\section{Introduction}

The target of this draft is not to claim \(\GRH\).  The target is to make the
current proof shape editable: every claimed step is written as a theorem or a
lemma, and every missing step is marked as a GAP with its current UV label.

The message is the following.  The RH manuscript appears to contain a local
phase, jet, whitening, and odd-projector package which is largely independent
of zeta once a source scalar has been supplied.  The first non-zeta source
target is the primitive Dirichlet pair
\[
  \Xi_\chi(s):=\Lambda(s,\chi)\Lambda(s,\overline\chi),
  \qquad
  \Phi_\chi^{\pair}(s):=\frac{\Xi_\chi(2s-1)}{\Xi_\chi(2s)}.
\]
For this object, the completed-zero source package can be written cleanly on
regular compact boundary intervals.  The remaining obstruction is not the
source sign or a hidden factor of \(2\).  It is the local normal-form question:
does the paired local block use the same scalar as its zeroth value
coordinate?

\subsection{What this draft proves}

The unconditional theorem in this draft is the completed paired source package
for primitive nonprincipal characters; see Theorem~\ref{thm:paired-source}.
It proves the sign-audited identity
\[
  q_\chi^{\pair}(t)
  =
  \Real\sum_{\rho\in Z(\Xi_\chi)}
  m_\rho
  \left(
    \frac{1}{1+2it-\rho}
    -
    \frac{1}{2it-\rho}
  \right)
\]
on regular compact intervals, with the completed-Hadamard normalization
\[
  B_{\chi,\comp}^{\pair}=0,
  \qquad
  S_{\chi,\comp}^{\pair}=q_\chi^{\pair}.
\]

\subsection{What this draft does not prove}

This draft does not prove \(\GRH\), \(\ERH\), or a tau-family analogue.  Those
statements require the source-to-slot, odd/transverse, remainder, and final
contradiction layers listed in Sections~\ref{sec:slot}--\ref{sec:transfer}.
They are recorded as GAPs rather than hidden in prose.

\section{The zeta-side input}
\label{sec:zeta-input}

The zeta case is treated as special in this paper.  The RH manuscript is the
place where the zeta source theorem, local slot theorem, odd transverse
package, and final contradiction are meant to be proved.  This draft only
records the interface required from that manuscript.

\begin{hypothesis}[Zeta source interface]
\label{hyp:zeta-source}
Let
\[
  \phi_\zeta(s)=\frac{\Lambda_\zeta(2s-1)}{\Lambda_\zeta(2s)}.
\]
On each regular compact boundary interval, the RH manuscript supplies a
phase convention, a source derivative \(q_\zeta\), a background term
\(B_\zeta\), and a positive strip-edge source scalar \(S_\zeta\) such that
\[
  q_\zeta(t)=B_\zeta(t)+S_\zeta(t).
\]
The source scalar is the same scalar that enters the RH local value slot.
\end{hypothesis}

\begin{gap}[UV-019: zeta source theorem]
\label{gap:zeta-source}
The canonical zeta source theorem for \(q_\zeta\), \(B_\zeta\), and
\(S_\zeta\) is delegated to the RH paper.  The GRH transfer draft should not
reprove it and should not use it without the phase convention, compact
summation, background naming, and multiplicity convention fixed there.
\end{gap}

\section{Primitive Dirichlet completed functions}
\label{sec:dirichlet-input}

Let \(\chi\) be a primitive nonprincipal Dirichlet character modulo \(r\), and
put
\[
  a_\chi:=\frac{1-\chi(-1)}{2}.
\]
We use the completed \(L\)-function normalization
\[
  \Lambda(s,\chi)
  :=
  \left(\frac{r}{\pi}\right)^{(s+a_\chi)/2}
  \Gamma\!\left(\frac{s+a_\chi}{2}\right)L(s,\chi).
\]

\begin{proposition}[Standard completed Dirichlet input]
\label{prop:standard-input}
For primitive nonprincipal \(\chi\), the completed function
\(\Lambda(s,\chi)\) is entire of order at most one and satisfies the standard
primitive functional equation relating \(\chi\) and \(\overline\chi\).  The
paired product
\[
  \Xi_\chi(s):=\Lambda(s,\chi)\Lambda(s,\overline\chi)
\]
is entire of order at most one.  Its zeros carry analytic multiplicities, and
the nonzero zeros satisfy
\[
  \sum_{\rho\neq 0}m_\rho |\rho|^{-2}<\infty.
\]
\end{proposition}

\begin{proof}
The continuation, primitive functional equation, trivial-zero locations, and
\(L(1,\chi)\neq 0\) input are the standard Dirichlet \(L\)-function facts,
for example in DLMF Section~25.15.  The displayed completed normalization is
obtained from the primitive functional equation by the parity choice
\(a_\chi=(1-\chi(-1))/2\).  The gamma poles are cancelled by the corresponding
trivial zeros, so the completed primitive nonprincipal function is entire.

For this working draft, the order-one input is taken in the standard bounded
form: combine the DLMF growth formulae for \(L(s,\chi)\), the primitive
functional equation, and Stirling's formula for \(\Gamma\).  This gives finite
order at most one for \(\Lambda(s,\chi)\), hence also for
\(\Xi_\chi\).  Hadamard factorization for entire functions of finite order then
gives a genus-one product, analytic zero multiplicities, and the summability
\(\sum m_\rho|\rho|^{-2}<\infty\).  A final publication draft should replace
this paragraph by exact textbook page or theorem references.
\end{proof}

\begin{gap}[Bibliography upgrade for Proposition~\ref{prop:standard-input}]
\label{gap:bibliography}
The working citation stack is DLMF Section~25.15, DLMF Section~5.11 for
Stirling, and a standard Hadamard factorization theorem such as Conway,
Chapter XI, Section~3.  A final draft should verify exact Davenport or
Montgomery--Vaughan page and equation references.
\end{gap}

\section{The completed paired source package}
\label{sec:source}

\begin{definition}[Paired quotient]
For primitive nonprincipal \(\chi\), define
\[
  \Phi_\chi^{\pair}(s)
  :=
  \frac{\Xi_\chi(2s-1)}{\Xi_\chi(2s)}.
\]
Let \(I\subset \R\) be compact and assume that
\[
  2it\notin Z(\Xi_\chi),
  \qquad
  1+2it\notin Z(\Xi_\chi)
\]
for every \(t\in I\).  Such an interval is called regular.
\end{definition}

\begin{definition}[Boundary phase and source derivative]
\label{def:boundary-phase}
On a regular compact interval \(I\), choose a real \(C^1\) boundary phase
\(\Theta_\chi^{\pair}\) by
\[
  \Phi_\chi^{\pair}\!\left(\frac12+it\right)
  =
  e^{-2i\Theta_\chi^{\pair}(t)}.
\]
Set
\[
  q_\chi^{\pair}(t):=
  \Theta_\chi^{\pair\,\prime}(t).
\]
\end{definition}

\begin{lemma}[Sign bridge]
\label{lem:sign-bridge}
On a regular compact interval,
\[
  q_\chi^{\pair}(t)
  =
  -\frac12
  \frac{\Phi_\chi^{\pair\,\prime}}{\Phi_\chi^{\pair}}
  \!\left(\frac12+it\right).
\]
Equivalently, if
\[
  D_\chi(t):=
  \frac{\Xi_\chi'}{\Xi_\chi}(1+2it)
  -
  \frac{\Xi_\chi'}{\Xi_\chi}(2it),
\]
then \(q_\chi^{\pair}(t)=D_\chi(t)\).  In particular \(D_\chi(t)\) is real on
the chosen boundary branch.
\end{lemma}

\begin{proof}
Differentiate
\[
  \Phi_\chi^{\pair}\!\left(\frac12+it\right)
  =
  e^{-2i\Theta_\chi^{\pair}(t)}
\]
with respect to \(t\).  Since \(s=1/2+it\), the left side contributes
\[
  i\frac{\Phi_\chi^{\pair\,\prime}}{\Phi_\chi^{\pair}}
  \!\left(\frac12+it\right),
\]
while the right side contributes \(-2i\Theta_\chi^{\pair\,\prime}(t)\).  Hence
\[
  \frac{\Phi_\chi^{\pair\,\prime}}{\Phi_\chi^{\pair}}
  \!\left(\frac12+it\right)
  =
  -2q_\chi^{\pair}(t).
\]
The chain rule for
\(\Phi_\chi^{\pair}(s)=\Xi_\chi(2s-1)/\Xi_\chi(2s)\) gives
\[
  \frac{\Phi_\chi^{\pair\,\prime}}{\Phi_\chi^{\pair}}(s)
  =
  2\frac{\Xi_\chi'}{\Xi_\chi}(2s-1)
  -
  2\frac{\Xi_\chi'}{\Xi_\chi}(2s).
\]
Putting \(s=1/2+it\) yields
\[
  -\frac12
  \frac{\Phi_\chi^{\pair\,\prime}}{\Phi_\chi^{\pair}}
  \!\left(\frac12+it\right)
  =
  \frac{\Xi_\chi'}{\Xi_\chi}(1+2it)
  -
  \frac{\Xi_\chi'}{\Xi_\chi}(2it)
  =
  D_\chi(t).
\]
\end{proof}

\begin{theorem}[Completed paired source package]
\label{thm:paired-source}
Let \(\chi\) be primitive nonprincipal, and let \(I\subset\R\) be a regular
compact interval.  Then
\[
  q_\chi^{\pair}(t)
  =
  \Real\sum_{\rho\in Z(\Xi_\chi)}
  m_\rho
  \left(
    \frac{1}{1+2it-\rho}
    -
    \frac{1}{2it-\rho}
  \right),
  \qquad t\in I,
\]
where the sum runs over distinct completed zeros of \(\Xi_\chi\), and
\(m_\rho\) is analytic multiplicity in
\(\Xi_\chi=\Lambda(\cdot,\chi)\Lambda(\cdot,\overline\chi)\).  The series
converges absolutely and uniformly on \(I\), and the same is true after any
fixed number of \(t\)-derivatives.

Thus, in completed-Hadamard normalization,
\[
  q_{\chi,\comp}^{\pair}(t)
  =
  B_{\chi,\comp}^{\pair}(t)
  +
  S_{\chi,\comp}^{\pair}(t),
  \qquad
  B_{\chi,\comp}^{\pair}(t)=0,
\]
with
\[
  S_{\chi,\comp}^{\pair}(t)
  :=
  \Real\sum_{\rho\in Z(\Xi_\chi)}
  m_\rho
  \left(
    \frac{1}{1+2it-\rho}
    -
    \frac{1}{2it-\rho}
  \right).
\]
\end{theorem}

\begin{proof}
By Proposition~\ref{prop:standard-input}, \(\Xi_\chi\) is entire of order at
most one.  On any simply connected region avoiding its zeros, the genus-one
Hadamard product gives the logarithmic derivative in the form
\[
  \frac{\Xi_\chi'}{\Xi_\chi}(w)
  =
  b_\chi+
  \sum_{\rho\in Z(\Xi_\chi)}
  m_\rho
  \left(
    \frac{1}{w-\rho}+\frac{1}{\rho}
  \right),
\]
with locally uniform convergence after the usual regularization.  If a zero at
\(\rho=0\) is admitted by a broader convention, it is separated as a finite
factor.  In the primitive nonprincipal setting used here, the boundary edge
points in \(I\) avoid zeros by regularity.

Take the edge difference at \(w=1+2it\) and \(w=2it\).  The constant
\(b_\chi\) cancels, and the regularizing \(1/\rho\) terms cancel term by term.
Hence
\[
  D_\chi(t)
  =
  \sum_{\rho\in Z(\Xi_\chi)}
  m_\rho
  \left(
    \frac{1}{1+2it-\rho}
    -
    \frac{1}{2it-\rho}
  \right).
\]
For \(t\in I\), regularity gives a positive distance from the two moving edge
points to the zero set.  For large \(|\rho|\), each summand is
\[
  O_I(|\rho|^{-2}),
\]
because it is the difference of two resolvents whose arguments differ by
\(1\).  The zero summability from Proposition~\ref{prop:standard-input} then
gives absolute and uniform convergence on \(I\).  After \(k\) derivatives in
\(t\), the summand is a constant multiple of the difference of two
\((k+1)\)-st powers of the same resolvents, hence is
\(O_{I,k}(|\rho|^{-k-2})\).  This proves uniform convergence after every fixed
number of derivatives.

By Lemma~\ref{lem:sign-bridge}, \(q_\chi^{\pair}=D_\chi\), and the boundary
phase makes this quantity real.  Keeping \(\Real\) in the display records the
scalar source identity.  Since the completed-zero theorem already sums every
completed zero of \(\Xi_\chi\) with analytic multiplicity, there is no further
completed background term.  This is exactly
\(B_{\chi,\comp}^{\pair}=0\).
\end{proof}

\begin{remark}[Raw \(L\)-factor bookkeeping]
The theorem uses completed zeros of
\(\Xi_\chi\).  In this normalization there is no additional gamma,
trivial-zero, pole, conductor, or Hadamard-exponential background term.  If one
rewrites the same identity using raw zeros of \(L(s,\chi)\) and
\(L(s,\overline\chi)\), those raw background pieces must be placed in a
separate raw bookkeeping term.  They must not be added to
Theorem~\ref{thm:paired-source}.
\end{remark}

\begin{gap}[UV-016: final promotion bookkeeping]
\label{gap:uv016}
The mathematical wording of Theorem~\ref{thm:paired-source} follows the
verifier-cleared staged source theorem.  The remaining UV-016 work for a final
paper is bibliographic polish and any decision about whether to keep the
completed normalization \(B_{\chi,\comp}^{\pair}=0\) as the main theorem or add
a separate raw-normalization corollary.
\end{gap}

\begin{gap}[UV-021: sign convention installed here]
\label{gap:uv021}
The sign convention is installed in Definition~\ref{def:boundary-phase} and
Lemma~\ref{lem:sign-bridge}: fixed quotient, negative boundary exponent,
\(q=\Theta'\), and
\[
  q_\chi^{\pair}
  =
  -\frac12
  (\Phi_\chi^{\pair\,\prime}/\Phi_\chi^{\pair})(1/2+it).
\]
Older positive-exponent text should be treated as obsolete unless the quotient
is inverted or \(q\) is redefined.
\end{gap}

\section{The source-to-slot layer}
\label{sec:slot}

The source package is not yet a \(\GRH\) theorem.  It must feed the same local
value channel used by the RH manuscript.  The following theorem records the
exact conditional mechanism.

\begin{hypothesis}[Paired unit-coordinate local chart]
\label{hyp:unit-coordinate}
For each \(m\in I\), the corrected paired local block is represented near the
background by a holomorphic map
\[
  \calF_{\chi,m,\sigma}(a,\eta),
\]
where \(a\) is the zeroth value coordinate
\[
  a=q_\chi^{\pair}(m)-B_\chi^{\pair}(m)
\]
and \(\eta\) contains all non-value local coordinates: derivative, mixed-jet,
curvature, background, cutoff, correction, multiplicity, and normalization
data.  In the completed normalization of Theorem~\ref{thm:paired-source},
\[
  a=S_{\chi,\comp}^{\pair}(m)=q_\chi^{\pair}(m).
\]
The pure value path is
\[
  (a,\eta)=(\alpha,\eta_0),
\]
so all non-value coordinates are frozen.
\end{hypothesis}

\begin{hypothesis}[Holomorphic whitening]
\label{hyp:whitening}
The corrected paired same-point and mixed blocks are holomorphic in
\((a,\eta)\) near the background, the same-point blocks are positive and
nondegenerate there, and the inverse-square-root whitening map is holomorphic.
\end{hypothesis}

\begin{theorem}[Conditional paired source-to-slot realization]
\label{thm:conditional-slot}
Assume Theorem~\ref{thm:paired-source},
Hypothesis~\ref{hyp:unit-coordinate}, and
Hypothesis~\ref{hyp:whitening}.  Define
\[
  A_{\val,\chi}^{\pair}(m,\sigma)
  :=
  \partial_a\calF_{\chi,m,\sigma}(0,\eta_0)
\]
and
\[
  R_\chi^{\pair}(m,\sigma)
  :=
  \calF_{\chi,m,\sigma}
  \bigl(S_{\chi,\comp}^{\pair}(m),\eta_\chi(m)\bigr)
  -
  \calF_{\chi,m,\sigma}(0,\eta_0)
  -
  S_{\chi,\comp}^{\pair}(m)
  A_{\val,\chi}^{\pair}(m,\sigma).
\]
Then
\[
  \Delta_\chi^{\pair}(m,\sigma)
  =
  S_{\chi,\comp}^{\pair}(m)
  A_{\val,\chi}^{\pair}(m,\sigma)
  +
  R_\chi^{\pair}(m,\sigma).
\]
Moreover, along the pure value path the remainder has zero first derivative
after subtracting the value term:
\[
  \left.
  \frac{\partial}{\partial \alpha}
  \left(
    \calF_{\chi,m,\sigma}(\alpha,\eta_0)
    -
    \calF_{\chi,m,\sigma}(0,\eta_0)
    -
    \alpha A_{\val,\chi}^{\pair}(m,\sigma)
  \right)
  \right|_{\alpha=0}
  =
  0.
\]
\end{theorem}

\begin{proof}
This is the local algebra of a holomorphic coordinate chart.  By
Hypothesis~\ref{hyp:unit-coordinate}, the source scalar from
Theorem~\ref{thm:paired-source} is the value coordinate \(a\).  By
Hypothesis~\ref{hyp:whitening}, the corrected whitened block is holomorphic in
that coordinate and in the frozen non-value coordinates.  The pure value
derivative at the background is therefore
\[
  \partial_a\calF_{\chi,m,\sigma}(0,\eta_0)
  =
  A_{\val,\chi}^{\pair}(m,\sigma).
\]
The displayed definition of \(R_\chi^{\pair}\) is then an identity.  The last
display follows by differentiating the same expression along
\((a,\eta)=(\alpha,\eta_0)\) at \(\alpha=0\).  The proof is non-tautological
only because the value coordinate was fixed before \(R_\chi^{\pair}\) was
defined.
\end{proof}

\begin{corollary}[No hidden scalar if the tested value derivative is nonzero]
\label{cor:no-hidden-scalar}
Assume the hypotheses of Theorem~\ref{thm:conditional-slot} and suppose the
tested value derivative \(A_{\val,\chi}^{\pair}(m,\sigma)\) is nonzero in the
slot under consideration.  A replacement
\[
  c_\chi(m)S_{\chi,\comp}^{\pair}(m)
\]
can satisfy the same pure value derivative test only if \(c_\chi(m)=1\).
\end{corollary}

\begin{proof}
If the coordinate were
\[
  a=c_\chi(m)S_{\chi,\comp}^{\pair}(m)+O(S^2),
\]
then the first derivative of the pure value path with respect to the source
scalar would be
\[
  c_\chi(m)A_{\val,\chi}^{\pair}(m,\sigma).
\]
Equality with the unit coefficient derivative from
Theorem~\ref{thm:conditional-slot} forces \(c_\chi(m)=1\) when the tested
derivative is nonzero.
\end{proof}

\begin{gap}[UV-017: paired unit-coordinate normal form]
\label{gap:uv017}
The missing proof is Hypothesis~\ref{hyp:unit-coordinate}.  One must construct
the paired corrected local blocks from the same boundary phase
\(\Theta_\chi^{\pair}\) and prove that their zeroth coordinate is exactly
\[
  q_\chi^{\pair}(m)-B_\chi^{\pair}(m),
\]
not \(c_\chi(m)(q-B)\) and not an upstairs-only surrogate.  The paired
holomorphic-whitening hypotheses and the scalar readout normalization also
belong to UV-017 unless the theorem is deliberately stated with them as
assumptions.
\end{gap}

\section{Odd/transverse and remainder layers}
\label{sec:odd}

The RH manuscript also uses a corrected odd transverse scalar, boundary
control, and an \(N\)-point odd-moment projector.  The paired Dirichlet source
package does not supply these ingredients.

\begin{hypothesis}[Odd/transverse package]
\label{hyp:odd}
For the family under consideration, the corrected local source-to-slot block
produces an odd holomorphic transverse germ with the same kind of boundary
control used in the RH manuscript.  The discrete \(N\)-point odd projector
annihilates the lower odd jets and leaves only the controlled high-order
transverse tail.
\end{hypothesis}

\begin{hypothesis}[Remainder dominance]
\label{hyp:remainder}
The local remainder \(R_\chi^{\pair}\) is smaller than the leading value-channel
term in the calibrated regime required by the RH contradiction argument.
\end{hypothesis}

\begin{gap}[UV-018: corrected odd/transverse realization]
\label{gap:uv018}
The corrected odd transverse scalar, boundary estimate, and \(N\)-point
projector input have not been constructed for primitive Dirichlet pairs.  This
is downstream from UV-017 and should not be folded into the source theorem.
\end{gap}

\begin{gap}[Remainder dominance]
\label{gap:remainder}
The current GRH team ledger does not yet close the paired remainder dominance
needed to compare the leading value term with the transverse/projected error
terms.  This is a separate analytic estimate, not a source-normalization
problem.
\end{gap}

\section{Formal transfer theorem}
\label{sec:transfer}

This section states the paper's long target in a form that can be revised as
gaps close.

\begin{hypothesis}[RH local contradiction template]
\label{hyp:rh-template}
The RH manuscript supplies the following local contradiction template.  If a
family has:
\begin{enumerate}
\item a source scalar occupying the exact local value slot;
\item holomorphic whitening and same-point positivity;
\item odd/transverse realization and boundary control;
\item remainder dominance in the calibrated regime;
\item the same off-line toy anomaly comparison;
\end{enumerate}
then an off-line zero in that family produces an incompatible transverse local
configuration.
\end{hypothesis}

\begin{theorem}[Conditional primitive Dirichlet GRH transfer]
\label{thm:conditional-grh}
Let \(\chi\) be primitive nonprincipal.  Assume the completed paired source
package of Theorem~\ref{thm:paired-source}, the paired source-to-slot
realization of Theorem~\ref{thm:conditional-slot}, the odd/transverse package
Hypothesis~\ref{hyp:odd}, the remainder dominance
Hypothesis~\ref{hyp:remainder}, and the RH local contradiction template
Hypothesis~\ref{hyp:rh-template}.  Then the paired local model admits no
off-line zero configuration for \(L(s,\chi)\) and \(L(s,\overline\chi)\) in the
scope covered by those hypotheses.
\end{theorem}

\begin{proof}
Theorem~\ref{thm:paired-source} supplies the completed source scalar with the
correct sign and normalization.  Theorem~\ref{thm:conditional-slot} identifies
that scalar with the local value slot.  Hypothesis~\ref{hyp:odd} supplies the
odd transverse channel and the finite projector.  Hypothesis~\ref{hyp:remainder}
makes the residual terms subordinate to the leading calibrated comparison.
Hypothesis~\ref{hyp:rh-template} then applies the RH manuscript's local
contradiction argument to the paired Dirichlet local model.  Therefore the
off-line configuration is excluded inside the stated scope.
\end{proof}

\begin{gap}[Not yet a proof of GRH]
\label{gap:not-grh}
Theorem~\ref{thm:conditional-grh} is conditional.  It becomes a GRH-family
theorem only after UV-017, UV-018, remainder dominance, and the RH contradiction
template are promoted without assumptions.
\end{gap}

\section{Tau and other families}
\label{sec:tau}

The Ramanujan tau \(L\)-function remains a useful stress test because it is
self-dual, but the current source and local realization package is weaker than
the paired Dirichlet package.

\begin{gap}[UV-020: tau compact source and localization package]
\label{gap:uv020}
For tau, one still needs a compact source decomposition, microscopic
denominator and holomorphy control, the correct block scaling hierarchy,
whitening gain and boundary estimates, and growth bookkeeping.  This draft
does not attempt to prove those statements.
\end{gap}

\section{Current proof ledger}
\label{sec:ledger}

\begin{itemize}
\item Completed paired source theorem: proved for the working draft in
Theorem~\ref{thm:paired-source}; final citation polish is
GAP~\ref{gap:bibliography}.
\item Paired sign convention: installed in Lemma~\ref{lem:sign-bridge};
GAP~\ref{gap:uv021} records the obsolete convention risk.
\item Zeta source theorem: delegated to the RH manuscript via
Hypothesis~\ref{hyp:zeta-source} and GAP~\ref{gap:zeta-source}.
\item Source-to-slot realization: conditional in
Theorem~\ref{thm:conditional-slot}; the missing local normal form is
GAP~\ref{gap:uv017}.
\item Odd/transverse package: missing, as recorded by
Hypothesis~\ref{hyp:odd} and GAP~\ref{gap:uv018}.
\item Remainder dominance: missing, as recorded by
Hypothesis~\ref{hyp:remainder} and GAP~\ref{gap:remainder}.
\item GRH-family conclusion: conditional only; see
Theorem~\ref{thm:conditional-grh} and GAP~\ref{gap:not-grh}.
\item Tau: exploratory; see GAP~\ref{gap:uv020}.
\end{itemize}

\section{Provenance}

This skeleton is based on the current GRH team directory under
\path{grh/teams/}, with timestamp \path{20260423-050630} and slug
\path{research-team-grh-dirichlet-tau}.
The core source theorem is drawn from \path{paper-updates.md}, the compact
regularization reports, the background-multiplicity reports, and the Sartre
verifier reports in the same team directory.  The UV-017 reduction is drawn
from Noether's coefficient-freeze report and its
\path{coefficient_freeze_reduction.md} note.  The live UV references are
UV-016 through UV-021 in \path{uv.md}.

\begin{thebibliography}{9}

\bibitem{DLMF25}
NIST Digital Library of Mathematical Functions, Section 25.15,
\emph{Dirichlet \(L\)-functions}.
\url{https://dlmf.nist.gov/25.15}

\bibitem{DLMF5}
NIST Digital Library of Mathematical Functions, Section 5.11,
\emph{Gamma function asymptotics}.
\url{https://dlmf.nist.gov/5.11}

\bibitem{Conway}
J. B. Conway,
\emph{Functions of One Complex Variable I}, second edition,
Springer, Chapter XI, Section 3.

\end{thebibliography}

\end{document}
